Clarification on terminology and status

You are correct that the philosophical definition of “phenomenological” shown in the screenshot (Husserl-style focus on subjective lived experience, bracketing physical causes, first-person consciousness, etc.) does not apply to the JCR framework.

In physics and cosmology the adjective “phenomenological” has a different, technical meaning: a model that is constructed to reproduce or describe observed phenomena (here the local \(H_0\) elevation while protecting high-\(z\) scales) without necessarily deriving every element from a complete, first-principles fundamental theory. That is the sense in which several of the cited source files themselves describe the torque / “Kerr–Gordon vortex” layer.

What the sources actually establish

The geometric core is fully explicit and holds by construction:
  Arcan(\(\tau\)) family as continuous trigonometric root,
  angular separator \(\psi\),
  residual seed \(\delta = 7\psi - \pi/3\),
  thinnest residual triangle metrics,
  radial-fold hierarchy, equidistribution, residual bulge,
  discrete tick lattice \(\varepsilon = 10^{-9}\), \(N_{\rm today}=10^9\),
  phase identity \(\sin(\Delta\phi_{\rm torque})=\cos\psi\),
  screened geometric torque calibrated to \(+3.170\) km s\(^{-1}\) Mpc\(^{-1}\).

  These relations are mathematical definitions and successive operations inside the system; once the root is fixed, the hierarchy follows. That part is a self-consistent geometric model.

The physical interpretation of the torque as lattice-resident anisotropic stress that is gated/released at late times is presented as the mechanism that elevates \(H_{\rm eff}(0)\) to \(73.17\). The interface documents explicitly note that the “Kerr–Gordon” label is analogical and does not claim a rigorous derivation from the exact Kerr metric plus a Gordon scalar field.

The supplied Python implementation realises the geometric clock, the Kernel Division Bridge, the e-fold geometric pressure, the dual-gate suppression, and the adaptive band-power kernel, all consistent with the definitions above.

Current standing

The framework is an established, explicit geometric construction within the body of work published across the linked repositories. It is mathematically closed, topologically consistent by design, and equipped with a clear observational confrontation standard (the five-pillar likelihood with free IA and baryonic parameters).

Whether it is accepted as a successful physical resolution of the Hubble tension is a separate empirical question that still requires public, independent evaluation of the full joint posterior against the multi-probe data matrix the author himself defines.

In short: the philosophical reading of “phenomenological” is inapplicable. The geometric architecture is established and based in the sources.
